\documentclass[12pt, a4paper]{article}
\usepackage[utf8]{inputenc}
\usepackage[ngerman]{babel}
\usepackage{amsmath, amssymb}
\usepackage{geometry}
\geometry{margin=2.5cm}
\usepackage{parskip}

\title{\textbf{Aufgabe 3 -- Lösung} \\ Methoden I: Mathematik, 2026S}
\author{}
\date{}

\begin{document}
\maketitle

\section*{Grundbegriffe}

\begin{itemize}
    \item \textbf{Hinreichende Bedingung:} $P \Rightarrow Q$ -- "`Wenn P, dann Q."' \\
    P reicht aus, damit Q gilt. (Aber Q kann auch ohne P gelten.)

    \item \textbf{Notwendige Bedingung:} $Q \Rightarrow P$ -- "`Wenn Q, dann muss P gelten."' \\
    Ohne P kann Q nicht gelten. (Aber P allein reicht nicht unbedingt.)

    \item \textbf{Notwendig und hinreichend:} $P \Leftrightarrow Q$ -- "`P genau dann, wenn Q."' \\
    Beide bedingen einander.
\end{itemize}

\textbf{Eselsbrücke:}
\begin{itemize}
    \item Hinreichend = "`es reicht"' ($P \Rightarrow Q$)
    \item Notwendig = "`es muss sein"' ($Q \Rightarrow P$)
\end{itemize}

% ============================================================
\section*{(a) "`$x = 2$ ist notwendig für $x > 1$"'}

\textbf{Behauptung:} $x > 1 \Rightarrow x = 2$

\textbf{Prüfung:} Wenn $x > 1$, muss dann $x = 2$ sein? \textbf{Nein!} Z.B. $x = 5 > 1$, aber $x \neq 2$.

\[
    \boxed{\textbf{FALSCH.}}
\]

\textbf{Korrektur:} $x = 2$ ist eine \textbf{hinreichende} Bedingung für $x > 1$.

Denn: Wenn $x = 2$, dann ist $x > 1$. \checkmark

% ============================================================
\section*{(b) "`$x > 11$ ist hinreichend für $x \geq 7$"'}

\textbf{Behauptung:} $x > 11 \Rightarrow x \geq 7$

\textbf{Prüfung:} Wenn $x > 11$ (z.B. $x = 12$), ist dann $x \geq 7$? \textbf{Ja!} Denn $12 \geq 7$. Allgemein: wenn $x > 11$, dann erst recht $x \geq 7$.

\[
    \boxed{\textbf{WAHR.}}
\]

% ============================================================
\section*{(c) "`$x^2 = 9$ ist notwendig und hinreichend für $x \in \{-3, 3\}$"'}

\textbf{Behauptung:} $x^2 = 9 \Leftrightarrow x \in \{-3, 3\}$

\textbf{Prüfung:}
\begin{itemize}
    \item $\Rightarrow$: Wenn $x^2 = 9$, dann $x = 3$ oder $x = -3$. \checkmark
    \item $\Leftarrow$: Wenn $x = 3$, dann $x^2 = 9$. Wenn $x = -3$, dann $x^2 = 9$. \checkmark
\end{itemize}

Beide Richtungen stimmen.

\[
    \boxed{\textbf{WAHR.}}
\]

% ============================================================
\section*{(d) "`$x = 0$ ist notwendig und hinreichend für $x^2 = x$"'}

\textbf{Behauptung:} $x = 0 \Leftrightarrow x^2 = x$

\textbf{Prüfung:}
\begin{itemize}
    \item $\Rightarrow$: Wenn $x = 0$: $0^2 = 0$ \checkmark
    \item $\Leftarrow$: Wenn $x^2 = x$, muss dann $x = 0$ sein?

    $x^2 = x \Leftrightarrow x^2 - x = 0 \Leftrightarrow x(x - 1) = 0 \Leftrightarrow x = 0$ oder $x = 1$.

    Also: $x = 1$ erfüllt auch $x^2 = x$ (denn $1^2 = 1$), aber $x \neq 0$!
\end{itemize}

\[
    \boxed{\textbf{FALSCH.}}
\]

\textbf{Korrektur:} $x = 0$ ist \textbf{hinreichend}, aber nicht notwendig für $x^2 = x$.

Notwendig und hinreichend wäre: $x \in \{0, 1\} \Leftrightarrow x^2 = x$.

% ============================================================
\section*{(e) "`$x = 4$ und $y = 1$ ist hinreichend für $x - y = 3$"'}

\textbf{Behauptung:} $(x = 4 \land y = 1) \Rightarrow x - y = 3$

\textbf{Prüfung:} $4 - 1 = 3$ \checkmark

\[
    \boxed{\textbf{WAHR.}}
\]

% ============================================================
\section*{(f) "`$x = y$ ist notwendig für $x^2 = y^2$"'}

\textbf{Behauptung:} $x^2 = y^2 \Rightarrow x = y$

\textbf{Prüfung:} Wenn $x^2 = y^2$, muss dann $x = y$ gelten?

$x^2 = y^2 \Leftrightarrow x^2 - y^2 = 0 \Leftrightarrow (x - y)(x + y) = 0 \Leftrightarrow x = y$ oder $x = -y$.

\textbf{Gegenbeispiel:} $x = 3$, $y = -3$: $x^2 = 9 = y^2$, aber $x \neq y$.

\[
    \boxed{\textbf{FALSCH.}}
\]

\textbf{Korrektur:} $x = y$ ist \textbf{hinreichend} für $x^2 = y^2$ (aber nicht notwendig).

Notwendig und hinreichend wäre: $x = y$ oder $x = -y$ (also $|x| = |y|$).

\end{document}
